% 60e Olympiade Internationale de Mathématiques — Bath, Royaume-Uni.
% Énoncés officiels en français (imo-official.org). Chaque problème vaut 7 points.

\begin{center}
  \emph{Premier jour --- mardi 16 juillet 2019 --- Durée : 4 heures 30}
\end{center}

\subsection*{Problème 1}
On note $\mathbb{Z}$ l'ensemble des entiers relatifs. Trouver toutes les fonctions
$f : \mathbb{Z} \to \mathbb{Z}$ telles que, pour tous les entiers $a$ et $b$,
\[
f(2a) + 2f(b) = f\bigl(f(a+b)\bigr).
\]

\subsection*{Problème 2}
Soit $A_1$ et $B_1$ deux points appartenant respectivement aux côtés $[BC]$ et $[AC]$ d'un triangle $ABC$.
Soit également $P$ et $Q$ deux points appartenant respectivement aux segments $[AA_1]$ et $[BB_1]$, de sorte
que les droites $(PQ)$ et $(AB)$ soient parallèles. Soit $P_1$ un point, situé sur la droite $(PB_1)$, tel que
$B_1$ se retrouve strictement entre $P$ et $P_1$, et tel que $\widehat{PP_1C} = \widehat{BAC}$. De même, soit
$Q_1$ un point, situé sur la droite $(QA_1)$, tel que $A_1$ se retrouve strictement entre $Q$ et $Q_1$, et tel
que $\widehat{CQ_1Q} = \widehat{CBA}$.

Démontrer que les points $P$, $Q$, $P_1$ et $Q_1$ sont cocycliques.

\subsection*{Problème 3}
Un réseau social compte $2019$ membres. Certains de ces membres sont amis l'un avec l'autre, la relation
d'amitié étant réciproque. Des événements du type décrit ci-dessous surviennent successivement, l'un après
l'autre :

Soit $A$, $B$ et $C$ trois membres tels que $A$ soit ami avec $B$ et $C$, mais que $B$ et $C$ ne soient pas
amis ; alors $B$ et $C$ deviennent amis, mais $A$ n'est plus ami ni avec $B$, ni avec $C$. Les autres
relations d'amitié entre membres ne changent pas durant cet événement.

Initialement, $1010$ membres ont $1009$ amis chacun, et $1009$ membres ont $1010$ amis chacun. Démontrer qu'il
existe une suite de tels événements à la suite desquels chaque membre aura au plus un ami.

\bigskip
\begin{center}
  \emph{Second jour --- mercredi 17 juillet 2019 --- Durée : 4 heures 30}
\end{center}

\subsection*{Problème 4}
Trouver tous les couples d'entiers naturels non nuls $(k, n)$ tels que
\[
k! = (2^n - 1)(2^n - 2)(2^n - 4)\cdots(2^n - 2^{n-1}).
\]

\subsection*{Problème 5}
La banque de Bath a émis des pièces dont une face est marquée de la lettre $H$ et l'autre face est marquée de
la lettre $T$. Morgane a aligné $n$ de ces pièces de gauche à droite. Elle réalise alors plusieurs fois de suite
l'opération suivante : si la lettre $H$ est visible sur exactement $k$ pièces, avec $k \geqslant 1$, alors
Morgane retourne la $k^{\text{ème}}$ pièce en partant de la gauche ; si $k = 0$, elle s'arrête. Par exemple,
si $n = 3$, le processus partant de la configuration $THT$ sera
\[
THT \to HHT \to HTT \to TTT\,;
\]
Morgane s'arrête donc au bout de $3$ opérations.
\begin{enumerate}
  \item[(a)] Démontrer que, quelle que soit la configuration initiale, Morgane doit s'arrêter au bout d'un
  nombre fini d'opérations.
  \item[(b)] Pour chaque configuration initiale $C$, on note $L(C)$ le nombre d'opérations que va réaliser
  Morgane avant de s'arrêter. Par exemple, $L(THT) = 3$ et $L(TTT) = 0$. Trouver la valeur moyenne des nombres
  $L(C)$ obtenus lorsque $C$ parcourt l'ensemble des $2^n$ configurations initiales possibles.
\end{enumerate}

\subsection*{Problème 6}
Soit $ABC$ un triangle dont tous les angles sont aigus et tel que $AB \neq AC$. On note $\omega$ le cercle
inscrit dans $ABC$, $I$ le centre de $\omega$, et $D$, $E$ et $F$ les points de tangence respectifs de $\omega$
avec les côtés $[BC]$, $[CA]$ et $[AB]$. Soit $R$ le point de $\omega$, autre que $D$, tel que la droite $(DR)$
soit perpendiculaire à $(EF)$. Soit $P$ le point d'intersection, autre que $R$, entre la droite $(AR)$ et le
cercle $\omega$. Enfin, soit $Q$ le point d'intersection, autre que $P$, entre les cercles circonscrits à
$PCE$ et à $PBF$.

Démontrer que les droites $(DI)$ et $(PQ)$ sont sécantes en un point appartenant à la perpendiculaire à $(AI)$
passant par $A$.
